\section{The full positive primitive of the proper tribonacci quotient}
Use the notation and actual polynomial division
$E=[2]^u[3]^r=(1+az+bz^2)J+R_0+R_1z$ of the registered
proper-quotient theorem, with $a=\lambda-1,b=1/\lambda$.
\begin{theorem}
For EVERY $u\ge4,r\ge2u$, the positive primitive
\[
       d_n=(J_n-J_{n-1})^+,\qquad J_{-1}=0,
\]
is globally LC, INCLUDING its rank-zero atom $d_0=J_0$.
The same holds in the entire registered norm cone
$r\ge18,u\ge0,b\sigma^u b^r\le1/64$.
\end{theorem}
Ordinary LC of J alone does not imply this assertion. The proof
uses an additional complete primitive entrance and the following
exact head corrections; neither is silently dropped.

\paragraph{Retain the complete constant and linear quotient corrections.}
For an old J with remainder (x,y), multiplication of E by [2]
changes the quotient to
\[
                         J'=[2]J+\lambda y.              \tag{1}
\]
For multiplication by [3] the exact quotient is
\[
 J'=[3]J+\lambda(x+y-a\lambda y)+\lambda y\,z.           \tag{2}
\]
For clarity, the complete remainders in these identities are
\[
 \begin{array}{ll}
 [2]:&(x-\lambda y)+(x+(1-\lambda a)y)z,\\
 [3]:&(x-C_0)+(x+y-aC_0-C_1)z,
 \end{array}
\]
where $C_0=\lambda(x+y-a\lambda y)$ and $C_1=\lambda y$.
Two scoped commutative-ring Lean reports verify the FULL
polynomial identities under $\lambda b=1$.

The registered nonsymmetric uniform-window closure says that
if a positive polynomial is LC and its entire positive primitive
is LC, both properties are preserved by [2] or [3] convolution.
It is used exactly with its FULL initial-atom hypothesis.
By (1), the actual new signed differences agree with those of
the convolution from index TWO onward. By (2), they agree from
index THREE onward. Their positive parts therefore agree at all
these indices as well, even across the later first-mode clipping.
Consequently all actual new primitive LC centres from FOUR onward
are inherited. A new head proof at centres ONE, TWO and THREE
suffices for both extensions. The full natural-index head/tail
assembly is also kernel checked.

\paragraph{All changed primitive head values and all three head minors.}
Use the ACTUAL target remainder, not the old remainder. Its norm
bound gives $|R_0|,|R_1|\le1/4$ throughout the norm cone and rich
family. Starting with
\[
 J_0=1-R_0,\quad J_1=w-aJ_0-R_1,\quad
 J_n=\sum_{j\le n/2}{r\choose j}{w-j\choose n-2j}
                       -aJ_{n-1}-bJ_{n-2}\ (2\le n\le4),
\]
put $d_0=J_0$ and $d_n=J_n-J_{n-1}$ for n=1..4.
All FIVE values and all THREE adjacent minors at centres 1..3
are strictly positive whenever w>=12, 0<=r<=w, and the whole
target remainder box holds. Thus these are the actual positive
primitive values; no clipping sign assumption is implicit.

The complete proof is a Bernstein certificate after
$w=12+X$, $r=\theta w$, $R_i=-1/4+\eta_i/2$.
Its five value lists have 2,5,11,13,21 coefficients, and its
three minor lists 28,56,99 coefficients: 235 in total.
Every coefficient is nonnegative and all at $X^0$ are positive.
The independent verifier proves ALL polynomial identities by
complete tensor interpolation with degree bound (6,3,2,2):
7*4*3*3=252 evaluations per identity, 2016 in total. This is an
identity certificate on an unbounded real cone, not a parameter test.

\paragraph{The entire norm-cone induction and the complete finite entrance.}
For the pure r=18 seed, the entire positive primitive INCLUDING
rank zero is directly checked. Ordinary quotient LC is already
proved by the previous theorem. Under further ternary extension
the old and new pure norms contract, so the old uniform-window
closure and the new head inequalities prove the primitive assertion
for every pure r>=18.

At fixed r>=18, append binary factors until reaching a target
$b\sigma^u b^r\le1/64$. All earlier norms are at most that target,
since sigma>1. At every step ordinary LC of the old and actual new
quotient is supplied by the registered norm-cone theorem. The old
primitive hypothesis comes from induction, the convolution primitive
is LC by the registered window closure, and (1) plus the complete
head proof give the target primitive. The same unbounded increasing-
norm predicate induction is reused; no new upper count bound occurs.

The remaining rich r<18 cases are exactly the previous THIRTY
explicit contexts. Their full primitives are now separately
reconstructed and certified, together with the pure r=18 seed.
There are 995 positive-part coefficient reads, including zeros,
and 933 LC minors, all nonnegative. Every nonzero signed primitive
coefficient is classified by a certified exact algebraic sign;
zero minors are exact field identities. The new finite certificate
contains 1693 distinct positive field triples in eighteen ordinary
Lean chunks. It is a NEW primitive certificate; the previous ordinary
quotient certificate is not mistaken for a primitive proof.

\paragraph{Scope.}
Four NEW reports comprise the two complete polynomial quotient
identities, the entire head/tail primitive join, and the aggregate
finite primitive/Bernstein certificate. Previous norm, cubic-root,
Bernstein-sum and nonsymmetric-window kernels are reused without
recounting. The actual quotient identifications, uniform-window
assembly, complete polynomial-identity certificates and entire
norm-cone induction are written proofs, not an end-to-end Lean
composition datatype formalization.
The theorem concerns J. The WHOLE actual infinite-ones residual
$\lambda d_{n-1}-d_n$ is a different sequence: even its positive
primitive fails LC before the D mode on the ACTUAL host
$1^{160}2^{20}3^{40}$, at shifted centre77<107.
The negative exact field minor is also covered by the aggregate
scalar certificate, and an independent positive suffix-box
reconstruction checks the finite host and the entire relevant reads.
Thus this theorem cannot be transferred
by silently identifying J with that residual. The complete actual
source and finite-initial-padding corrections still need treatment.
Actual u>=33 strong diagonal, general C/full A2 and ORIGINAL
undilated coefficient unimodality remain OPEN.
